\section*{الفصل \ref{ch:b2:ore-deposits} --- مكامن الخامات}

\begin{solution}{exo:b2:ore-deposits:1}
(أ) $\qty{0.5}{\percent} = \qty{5000}{ppm}$، و$5000/28 = 179$. (ب) $\qty{6}{g/t} = \qty{6000}{ng/g}$،
و$6000/1.5 = 4000$. (ج) يحوي الكروميت كسرًا كتليًا من الكروم قدره $2 \times 52.00/(55.85 + 104.00 + 64.00) = 0.4646$؛
ويحوي الكروميتيت $0.9 \times 0.4646 = 0.418$، أي \qty{418000}{ppm}،
و$418000/92 = 4540$.
\end{solution}

\begin{solution}{exo:b2:ore-deposits:2}
(أ) $914 \times 0.0053 = \qty{4.84}{Mt}$ من النحاس و$914 \times 10^6 \times 0.40 = \num{3.66e8}$ غرام،
أي \qty{366}{t} من الذهب؛ و$1640 \times 0.0044 = \qty{7.22}{Mt}$ و\qty{738}{t}. (ب) $1640 - 914 =
\qty{726}{Mt}$ تحوي $7.216 - 4.844 = \qty{2.37}{Mt}$ من النحاس، أي بعيار متوسّط
$2.372/726 = \qty{0.327}{\percent}$: فالكتل الإضافية تقع دون \omterm{def:b2:ore-deposits:grade}{عيار القطع} \qty{0.3}{\percent}
وفوقه بقليل. (ج) إنّهما يجيبان عن السؤال عند عياري قطع مختلفين؛ ولا أحد منهما هو حجم
المكمن بمعزل عن عيار قطعه.
\end{solution}

\begin{solution}{exo:b2:ore-deposits:3}
(أ) $(19.3 - 1.00)/(2.65 - 1.00) = 18.3/1.65 = 11.1$. (ب) تساوي السرعتين يقتضي $(\rho_g - 1)d_g^2 =
(\rho_q - 1)d_q^2$، ومن ثمّ $d_g/d_q = \sqrt{1.65/18.3} = 0.300$: فحبيبة الذهب تترسّب كحبيبة كوارتز
أكبر منها بنحو ثلاث مرّات. (ج) $d_q = 0.1\sqrt{3.7/1.65} = \qty{0.150}{mm}$.
\end{solution}

\begin{solution}{exo:b2:ore-deposits:4}
لكلّ وحدة مساحة وعند الكثافة نفسها، يحتفظ الغطاء بنحاسه ويتلقّى نحاس العمود
المغسول: $g = g_0(h_e + \eta h_l)/h_e$. (أ) $0.4 \times (50 + 150)/50 = \qty{1.6}{\percent}$.
(ب) $0.4 \times (50 + 0.8 \times 150)/50 = \qty{1.36}{\percent}$. (ج) $h_l = 50 \times (1.7/0.4 - 1)
= \qty{163}{m}$. وتقع هذه العيارات في المدى 0.4--\qty{1.7}{\percent} لأغطية الإثراء في المكامن
البورفيرية.
\end{solution}

\begin{solution}{exo:b2:ore-deposits:5}
(أ) البلاتين، $C_S = 10 \times 10^4 (R + 1)/(R + 10^4)$ بوحدة \unit{ppb}: \qty{1.00}{ppm} عندما $R = 100$،
و\qty{9.10}{ppm} عندما \num{1000}، و\qty{50.0}{ppm} عندما \num{10000}، و\qty{90.9}{ppm} عندما \num{100000}.
والنيكل، $C_S = 200 \times 500 (R + 1)/(R + 500)$ بوحدة ppm: \qty{1.68}{\percent}، و\qty{6.67}{\percent}،
و\qty{9.52}{\percent}، و\qty{9.95}{\percent}. (ب) $D\,C_L$: \qty{100}{ppm} من البلاتين،
و\qty{10}{\percent} من النيكل. (ج) يبقى النيكل في حدود \qty{5}{\percent} من سقفه ابتداءً من
$R = \num{10000}$، فلا يكاد محتواه يتغيّر على مدى عامل عشرة في $R$؛ أمّا البلاتين فلا يزال
يتضاعف تقريبًا هناك، فمحتواه يقيس $R$.
\end{solution}

\begin{solution}{exo:b2:ore-deposits:6}
(أ) $\mu_M = 2.373 - 0.364 - 2 = 0.009$؛ $\sigma_M = 0.702$؛ والوسيط $10^{0.009} = \qty{1.02}{Mt}$.
(ب) $z = (1 - 0.009)/0.702 = 1.412$، و$1 - \Phi(1.412) = 0.0790$: أي نحو 33 مكمنًا من 422.
(ج) $\sigma = \sqrt{0.654^2 + 0.199^2} = 0.684$، و$z = 1.449$، و$1 - \Phi = 0.0737$، أي نحو 31: فالارتباط
الضعيف يضيف بضعة مكامن عملاقة.
\end{solution}

\begin{solution}{exo:b2:ore-deposits:7}
متوسّط الكتل كلّها $0.40\,e^{0.18} = \qty{0.479}{\percent}$. عندما $c = 0.3$: $z = \ln(0.40/0.3)/0.60 =
0.4795$، و$\Phi(z) = 0.6842$، و$\Phi(z + s) = 0.8598$: $T = \qty{684}{Mt}$، و$\bar g = 0.479 \times
0.8598/0.6842 = \qty{0.602}{\percent}$، و$M = 1000 \times 0.00479 \times 0.8598 = \qty{4.12}{Mt}$.
وعندما $c = 0.6$: $z = -0.6758$، و$\Phi = 0.2496$، و$\Phi(z + s) = 0.4698$: أي \qty{250}{Mt} بعيار
\qty{0.901}{\percent}، و\qty{2.25}{Mt}. وعلى الشكل، عند عيار قطع \qty{0.3}{\percent} يقرأ المنحنى
الأزرق نحو \qty{680}{Mt} والأحمر نحو \qty{0.6}{\percent}.
\end{solution}

\begin{solution}{exo:b2:ore-deposits:8}
(أ) المائع: $\num{5e6}/(\num{1e-3} \times 0.5) = \qty{1.00e10}{t}$؛ \omterm{def:b2:magma-emplacement:magma}{والصهارة}: $\num{1e10}/0.035 =
\qty{2.86e11}{t}$. (ب) $\num{2.86e14}/2600 = \qty{1.10e11}{m^3}$، أي \qty{110}{km^3}، وهو مكعّب من \omterm{def:b2:magma-emplacement:magma}{الصهارة}
طول ضلعه نحو \qty{4.8}{km}. (ج) كانت \omterm{def:b2:magma-emplacement:magma}{الصهارة} تحوي $\num{2.86e11}
\times \num{50e-6} = \qty{1.43e7}{t}$ من النحاس، وحاز الخام $5/14.3 = 0.350$ منه.
\end{solution}

\begin{solution}{exo:b2:ore-deposits:9}
(أ) متوسّط اللوغاريتمات 1.411، وانحرافها المعياري 0.0664؛ والوسيط $10^{1.411} = \qty{25.8}{ppm}$؛
والعتبة $10^{1.411 + 0.133} = 10^{1.544} = \qty{35.0}{ppm}$. (ب) $(4 \times 400 + 36 \times 25.8)/40 =
\qty{63.2}{ppm}$، وهي فوق \qty{35.0}{ppm}: فالرواسب لا تزال شاذّة بعد تخفيفها بتسعة أجزاء من الرواسب
النظيفة. (ج) $4 \times 400 + (A - 4) \times 25.8 = 35.0\,A$ تعطي $A = 4 \times (400 -
25.8)/(35.0 - 25.8) = \qty{163}{km^2}$.
\end{solution}

\begin{solution}{exo:b2:ore-deposits:10}
(أ) $e^{su}\,e^{-u^2/2} = e^{-(u - s)^2/2 + s^2/2}$؛ ثم اقسم على $\sqrt{2\pi}$. (ب) $E[e^{\mu + su}] =
e^{\mu}\int e^{su}\varphi(u)\,\mathrm{d}u = e^{\mu + s^2/2}\int\varphi(u - s)\,\mathrm{d}u = e^{\mu +
s^2/2}$. (ج) $\Phi$ متزايدة و$s > 0$، ومن ثمّ $\Phi(z + s) > \Phi(z)$: فالكتل المحتفَظ بها أغنى من
الكتل المطروحة، ولذا يحتفظ كلّ عيار قطع بحصّة من الفلز أكبر من حصّته من الصخر؛ ويزداد
الفرق مع $s$، أي مع تغيّرية المكمن.
\end{solution}

\begin{solution}{exo:b2:ore-deposits:11}
(أ) $D(R + 1)/(R + D) = fD$ تعطي $R + 1 = fR + fD$، و$R = (fD - 1)/(1 - f)$. (ب) $f = 0.5$: $R = D - 2$،
أي 298 و\num{9998}؛ و$f = 0.9$: $R = 9D - 10$، أي \num{2690} و\num{89990}. (ج) يبلغ النيكل
\qty{50}{\percent} من سقفه داخل مدى خامات النيكل والنحاس (من 100 إلى \num{2000})، و90 فوقه بقليل؛
أمّا البلاتين فيحتاج إلى $R$ من \num{10000} إلى \num{100000}، وهو مدى حيود عناصر مجموعة البلاتين.
\end{solution}

\begin{solution}{exo:b2:ore-deposits:12}
(أ) مع $\ln M = m + s u$، تحقّق المكامن التي تفوق الكمّية $u > \Phi^{-1}(q)$؛ وفلزّها
$\int_{\Phi^{-1}(q)}^{\infty} e^{m + su}\varphi(u)\,\mathrm{d}u = e^{m + s^2/2}[1 - \Phi(\Phi^{-1}(q) -
s)]$، والقسمة على المجموع $e^{m + s^2/2}$ تعطي $\Phi(s - \Phi^{-1}(q))$. (ب) $s = 0.702 \ln 10 =
1.616$؛ $\Phi(1.616 - 1.282) = \Phi(0.334) = 0.631$. (ج) $\Phi(1.616 - 2.326) = \Phi(-0.710) = 0.239$،
مقابل 0.266 للمكامن الأربعة الكبرى: فأكبر المكامن أكبر قليلًا ممّا يتوقّعه القانون
اللوغاريتمي الطبيعي.
\end{solution}

\begin{solution}{pb:b2:ore-deposits:1}
\textbf{1.} \qty{330}{m}؛ $\sum l\,g = 138.4$، والمتوسّط $138.4/330 = \qty{0.419}{\percent}$.
\textbf{2.} $3.24/8 = \qty{0.405}{\percent}$؛ فكلّ فترة تمثّل طولًا من اللباب، ومن ثمّ
حجمًا من الصخر: ويجب أن تُحسب الفترة الطويلة أكثر من القصيرة.
\textbf{3.} $230/330 = 0.697$، بعيار $118.6/230 = \qty{0.516}{\percent}$.
\textbf{4.} $4190/28 = 150$.
\textbf{5.} $8100/28 = 289$.
\textbf{6.} $0.38\,e^{0.55^2/2} = 0.38\,e^{0.1513} = \qty{0.442}{\percent}$؛ والثقب، بعيار
\qty{0.419}{\percent}، يقع في حدود \qty{6}{\percent} منه: فهو ممثِّل.
\textbf{7.} $\pi \times 450^2 \times 700 = \qty{4.45e8}{m^3}$؛ $\times 2.7 = \num{1.20e9}$ طنّ،
أي \qty{1200}{Mt}.
\textbf{8.} $z = \ln(0.38/0.3)/0.55 = 0.4298$؛ $\Phi(z) = 0.6663$.
\textbf{9.} $z = 1.167$ و0.4298 و$-0.4990$؛ $\Phi = 0.8784$ و0.6663 و0.3089؛ $T = \qty{1050}{Mt}$
و\qty{800}{Mt} و\qty{371}{Mt}.
\textbf{10.} $\Phi(z + s) = \Phi(0.9798) = 0.8364$؛ $\bar g = 0.442 \times 0.8364/0.6663 =
\qty{0.555}{\percent}$.
\textbf{11.} $0.442 \times 0.9570/0.8784 = \qty{0.482}{\percent}$؛ $0.442 \times 0.5203/0.3089 =
\qty{0.744}{\percent}$.
\textbf{12.} الكتلة الطنّية $\times\, 370.7/799.6 = 0.464$؛ والعيار $\times\, 0.744/0.555 = 1.34$.
\textbf{13.} عند الوسيط، \qty{0.38}{\percent} ($z = 0$): \qty{600}{Mt} بعيار $0.442 \times
\Phi(0.55)/0.5 = 0.442 \times 0.7088/0.5 = \qty{0.627}{\percent}$.
\textbf{14.} $1200 \times 0.00442 = \qty{5.30}{Mt}$.
\textbf{15.} $(\log_{10} 1200 - 2.373)/0.654 = (3.0792 - 2.373)/0.654 = 1.080$؛ $\Phi = 0.8599$: أي
المئين 86، فهو مكمن كبير.
\textbf{16.} $(\log_{10} 0.442 + 0.364)/0.199 = (-0.3546 + 0.364)/0.199 = 0.047$؛ $\Phi = 0.5187$: عيار
عادي، في المئين 52.
\textbf{17.} 0.702 مع الارتباط، و0.684 من دونه.
\textbf{18.} $\mu_M = 0.009$؛ $(\log_{10} 5.30 - 0.009)/0.702 = (0.7243 - 0.009)/0.702 = 1.019$؛ $\Phi =
0.8459$: أي المئين 85.
\textbf{19.} $422 \times (1 - 0.8459) = 65$.
\textbf{20.} $M = T_0\,e^{\mu + s^2/2}\,\Phi(z + s) = 1200 \times 0.00442 \times 0.9570 =
\qty{5.08}{Mt}$.
\textbf{21.} $1200 \times 0.00442 \times 0.5203 = \qty{2.76}{Mt}$.
\textbf{22.} $\Phi(z + s) = 0.8364$ من النحاس، في $\Phi(z) = 0.6663$ من الكتلة الطنّية.
\textbf{23.} من 0.3 إلى \qty{0.2}{\percent} تزداد الكتلة الطنّية بمقدار $1054 - 800 = \qty{254}{Mt}$
(\qty{32}{\percent}) والنحاس بمقدار $5.08 - 4.44 = \qty{0.64}{Mt}$ (\qty{14}{\percent}): فمتوسّط عيار الكتل
المضافة $0.64/254 = \qty{0.25}{\percent}$، وهي أفقر كتل المكمن.
\textbf{24.} $1200 \times 0.00442 \times 0.8364 = 800 \times 0.00555$: يحوي الموقع الواعد
\textbf{\qty[text-series-to-math=true]{4.44}{Mt}} من النحاس فوق \omterm{def:b2:ore-deposits:grade}{عيار القطع} \qty{0.3}{\percent}.
\end{solution}
